\documentclass[12pt]{article}
\usepackage[utf8]{inputenc}
\usepackage[T1]{fontenc}
\usepackage[brazilian]{babel}
\usepackage{geometry}
\usepackage{amsmath}
\usepackage{listings}
\usepackage{xcolor}
\usepackage{fancyhdr}
\usepackage{titlesec}
\usepackage{enumitem}

\geometry{margin=1in}
\setlength{\headheight}{15pt}
\setlength{\parindent}{0pt}
\setlength{\parskip}{4pt}

\definecolor{primary}{HTML}{1F4E79}
\definecolor{secondary}{HTML}{0B6E99}
\definecolor{lightbg}{HTML}{F4F8FC}
\definecolor{codebg}{HTML}{F7F7F7}
\definecolor{codeframe}{HTML}{D3D3D3}

\titleformat{\section}
  {\Large\bfseries\color{primary}}
  {}
  {0pt}
  {}

\titleformat{\subsection}
  {\large\bfseries\color{secondary}}
  {}
  {0pt}
  {}

\setlist[itemize]{leftmargin=1.4em,itemsep=2pt,topsep=3pt}
\setlist[enumerate]{leftmargin=1.6em,itemsep=2pt,topsep=3pt}

\pagestyle{fancy}
\fancyhf{}
\lhead{\textbf{Programação em Python}}
% \chead{Aula 02: Entrada e Conversão}
\rhead{Prof. Msc. Nator Junior Carvalho da Costa}
\cfoot{\thepage}
\renewcommand{\headrulewidth}{0.4pt}

\lstset{
    language=Python,
    basicstyle=\ttfamily\small,
    backgroundcolor=\color{codebg},
    frame=single,
    rulecolor=\color{codeframe},
    keywordstyle=\color{blue},
    commentstyle=\color{gray},
    stringstyle=\color{red},
    breaklines=true,
    showstringspaces=false,
    tabsize=4
}

\begin{document}

\begin{center}
{\LARGE \textbf{Lista de Atividades - Aula 02}}\\[4pt]
{\large Entrada de Texto e Conversão de Tipos}\\[10pt]
\textbf{Disciplina:} Lógica de Programação em Python\\
\textbf{Professor:} Prof. Msc. Nator Junior Carvalho da Costa\\
\textbf{Semestre:} 2026.1
\end{center}

\vspace{8pt}
\colorbox{lightbg}{%
\parbox{\dimexpr\linewidth-2\fboxsep\relax}{%
\textbf{Instruções}\\
Resolva as atividades abaixo usando Python. Lembre-se:
\begin{itemize}
    \item A função \texttt{input()} sempre retorna \texttt{str}.
    \item Converta os dados quando necessário usando \texttt{int()}, \texttt{float()} e \texttt{str()}.
    \item Use nomes claros para variáveis e exiba mensagens amigáveis ao usuário.
\end{itemize}
}}

\section*{Atividades}

\subsection*{Atividade 1: Cadastro Rápido}
Escreva um programa que:
\begin{enumerate}
    \item Solicite ao usuário o \texttt{nome} e a \texttt{idade}
    \item Converta a idade para inteiro com \texttt{int()}
    \item Exiba: \texttt{Olá, <nome>! Você tem <idade> anos.}
    \item Exiba também o tipo da variável idade usando \texttt{type()}
\end{enumerate}

\subsection*{Atividade 2: Soma de Dois Números}
Escreva um programa que:
\begin{enumerate}
    \item Solicite dois números reais ao usuário
    \item Converta ambos com \texttt{float()}
    \item Calcule a soma e armazene em \texttt{resultado}
    \item Exiba a soma com duas casas decimais
\end{enumerate}

\subsection*{Atividade 3: Conversor de Centímetros para Metros}
Escreva um programa que:
\begin{enumerate}
    \item Solicite um valor em centímetros (entrada do usuário)
    \item Converta o valor para \texttt{float}
    \item Calcule o valor equivalente em metros usando: $m = \frac{cm}{100}$
    \item Exiba no formato: \texttt{150 cm correspondem a 1.50 m}
\end{enumerate}

\subsection*{Atividade 4: Cálculo de Média com Entrada do Usuário}
Escreva um programa que:
\begin{enumerate}
    \item Solicite três notas do aluno pelo teclado
    \item Converta cada nota para \texttt{float}
    \item Calcule a média aritmética
    \item Exiba: \texttt{A média final é X.XX}
\end{enumerate}

\section*{Exemplo Base}
\begin{lstlisting}
nome = input("Digite seu nome: ")
idade = int(input("Digite sua idade: "))

print(f"Ola, {nome}! Voce tem {idade} anos.")
print(type(nome), type(idade))
\end{lstlisting}

\end{document}
